\section{DISCUSSION, RESEARCH ROADMAP, AND LIMITATIONS}
\label{sec:discussion_roadmap}

The preceding sections trace O-ISAC from physical sharing to metric evidence,
validation, application requirements, and network relevance. This section
brings those findings together, turns recurring evidence gaps into research
priorities, and states the boundaries within which the synthesis can be read.

\subsection{Interpretation Across the Evidence Base}

Physical integration identifies where communication and sensing share
infrastructure, signals, or information. Evidential integration shows how
clearly both outcomes remain traceable to the same operating conditions.
Together, these views connect syntheses of architecture and transmission
medium while preserving the physical regime of each platform
\cite{Wen2024OISACArchitectures,Su2025OISACTheoryPractice,
Liang2024LiSACReview,Liu2026OpticalFiberISACReview}.

Integration depth determines how many design variables become connected.
Sharing additional components can improve coordination across the two
functions. It also joins power budgets, distortion, synchronization, dynamic
range, and failure modes. The scientific value of deeper sharing therefore
depends on showing where these effects enter the signal chain and how they
shape both outcomes.

Tradeoff patterns transfer across studies when tasks, metric definitions,
measurement planes, baselines, constraints, and validation settings remain
aligned. This alignment allows model bounds and empirical operating points to
retain their distinct roles and support comparison within a stated context
\cite{Khorasgani2026GeneralOpticalISACSurvey,Rode2025PolarizationISAC}.

Fielded cable and network studies illustrate the value of operational
continuity
\cite{Cao2026SubmarineFiberISAC,Ip2026DeployedTelecomCableISAC}. Within the
survey, system evidence becomes reusable when configuration, calibration,
processing, and reported artifacts remain connected across interfaces. This
continuity links physical sharing to application and network claims. It also
defines the central research task, which is to establish joint operating
points that another group can interpret, test, and reuse.

\subsection{Research Priorities for Reusable Evidence}

Table~\ref{tab:research_roadmap} organizes the synthesis into five priorities
spanning measurement context, modality benchmarks, joint tests,
reconstruction, intelligence, scale, and security. Future studies provide the
evaluation setting.

\begin{table}[!t]
\caption{Research priorities derived from the reviewed evidence}
\label{tab:research_roadmap}
\centering
\footnotesize
\setlength{\tabcolsep}{2.4pt}
\renewcommand{\arraystretch}{1.02}
\begin{tabularx}{\columnwidth}{@{}
>{\raggedright\arraybackslash}p{0.20\columnwidth}
>{\raggedright\arraybackslash}X
>{\raggedright\arraybackslash}p{0.30\columnwidth}@{}}
\toprule
Priority & Finding and proposed action & Evidence of progress \\
\midrule
\textbf{1. Measurement contract} &
\emph{Finding.} Of 4,779 metric records, 118 supported bounded relations under
aligned conditions. \emph{Action.} Record the task, plane, unit, baseline,
operating condition, validation setting, and source locator. &
An independent reader reconstructs the quantity and condition from a machine
readable record. \\
\addlinespace[1.5pt]
\textbf{2. Modality specific benchmark} &
\emph{Finding.} Cross study comparison depends on matched physical context.
\emph{Action.} Repeat one task with reference and independent designs under
shared calibration and ground truth. &
Independent groups reproduce the relative finding under the shared protocol. \\
\addlinespace[1.5pt]
\textbf{3. Joint operation under disturbance} &
\emph{Finding.} Of 402 substantive tradeoff relationships, 371 depend on
reported conditions. Six of 12 field or deployment tier studies report
outcomes in both domains, while shared timing remains unresolved.
\emph{Action.} Evaluate both functions in the same run and condition set under
load, a modality specific disturbance, and matched single function baselines. &
Repeated trials recover the joint region and its limiting mechanism. \\
\addlinespace[1.5pt]
\textbf{4. Reconstructable experiment package} &
\emph{Finding.} Open data appear in 13 of 206 studies, and code or model access
reaches a partial release in one. \emph{Action.} Preserve a versioned package
and rerun the main comparison under one added condition. &
An independent rerun recovers the comparison and added condition from the
package. \\
\addlinespace[1.5pt]
\textbf{5. Intelligence, scale, and security} &
\emph{Finding.} Technology and application counts are multilabel. The
exclusive 6G crosswalk records corpus relevance, with conformance,
interoperability, and readiness kept as separate targets. \emph{Action.} Apply
a declared shift, attack, load, or component failure and report overhead,
resource use, decisions, and fallback. &
The service retains an explainable boundary with resource use, decisions, and
fallback recorded. \\
\bottomrule
\end{tabularx}
\vspace{1pt}
\parbox{\columnwidth}{\raggedright\footnotesize\emph{Note.} The findings
summarize the retained corpus presented in Sections~\ref{sec:optical_modality_families}
through \ref{sec:technologies_applications_6g}. The authors developed the
proposed actions and indicators of progress for future evaluation. The
numbered stages form a dependency order and retain the physical conditions of
each modality.\par}
\end{table}

The order matters because each stage uses evidence produced by the previous
one. Clear measurement context supports fair benchmarks. Benchmarks support
realistic joint tests. Reconstructable artifacts make those tests repeatable,
and repeatable tests provide a basis for adaptive, secure, and network scale
evaluation.

Modality specific stress tests should follow the native signal path. Fiber
tests can vary traffic and topology. Free space optical tests can vary
geometry, pointing, atmosphere, and blockage. Visible light tests can vary
illumination, field of view, orientation, and ambient light. Photonic terahertz
tests can vary optical front end drift together with wireless propagation.
Matched single function baselines keep the limiting mechanism visible across
these settings.

For standards and 6G coordination, this sequence provides a practical basis for
reporting profiles. A shared core can define essential metadata while
preserving modality physics and ground truth. Funding calls can support the
same goal through reusable condition descriptions, benchmark definitions, and
versioned artifacts.

\FloatBarrier
\subsection{Scope and Limitations of the Review}

The review searched six sources for reports published from 1 January 2020 to
22 June 2026 and required sufficient full technical content in English. Full
text assessment covered 272 of the 330 reports sought, with 58 recorded as
unretrieved. Exact query to export mapping remains incomplete for two low yield
Taylor \& Francis exports. These boundaries define the literature represented
by the 206 retained studies.

The OSF record was created retrospectively after searching and screening and
captures a 221 study predecessor state. Lineage reconciliation rebaselined that
predecessor state to the final 206 study corpus, whose attrition path follows
the executed search records in Section~\ref{sec:review_methods}. One
investigator supervised the workflow and made the final review decisions.
Routine independent duplicate screening, extraction, appraisal, and
adjudication were outside the executed workflow. Independent checks focused on
governed artifacts and reconciliation.

TQAF is a nonvalidated technical appraisal developed for this survey and
describes evidential support within the review. Conventional study level risk
of bias instruments and GRADE certainty assessment were outside the executed
appraisal. Task, quantity, and validation heterogeneity made structured
narrative synthesis the appropriate analytical form, with meta analysis
outside the review design. Formal sensitivity analysis requires a pooled or
model dependent primary effect beyond this structure. Formal missing results
and publication bias assessment requires harmonized prespecified outcomes and
source protocols, which were unavailable across the corpus. Selective
reporting and publication bias therefore remain possible. Causal moderator
inference belongs to a different analytical design.

Taxonomy and multilabel coding depend on reviewer judgment and can carry
classification error, especially for hybrid architectures. Row counts reflect
extraction granularity and represent coded observations. Independent effect
counts require a different unit of analysis. Ninety two legacy metric rows
received the conservative comparison status because the required conditions
were absent from their recorded fields. Artifact status reflects source
reporting at the audit date. Link verification was selective, and common
protocol execution of released artifacts remained outside the review.

The conclusions describe recurring mechanisms, measurement constraints, and
validation patterns within the retained corpus. Worldwide activity, market
share, deployment prevalence, causal effects, and platform ranking remain
separate research questions. Within these boundaries, the corpus supports a
clear conclusion about the conditions that turn optical integration into
transferable joint evidence. Section~\ref{sec:conclusion} brings these
conditions together in the final synthesis.
